% Appendix D, sections D.1--D.7: hand-typeset from monograph/registered-information-appendix-D.md
% (the Markdown source writes these sections' mathematics in plain text).
\appsection{Supplementary finite theorems}

This appendix is part of the monograph. D.1--D.7 integrate the results developed during the proof review. D.8--D.13 close additional finite questions about sites, witness certificates, propagation, robust transfer, the relationship to contextuality, and commutation of extraction. These are additions to this exposition, without a literature-priority claim.

All state and observation spaces are finite and nonempty. Information order is coarser-to-finer for partitions; $E\preceq_B F$ means that $E$ is a garbling of $F$, and $\delta(E,F)$ measures the error of simulating $F$ from $E$. A prior is specified for each Shannon quantity. When a hypothesis restricts the alphabet, coupling, site, or adversary, it is part of the theorem.

\subsecn{1}{D.1}{Maximal certified partial answer}

\begin{thmplain}[\hypertarget{stmt:D.1}{Theorem~D.1}]
Let a deterministic registered representation be $r:D\to X$ and a question be $q:D\to A$. With a fresh abstention symbol, a sound partial decoder $d:X\to A\cup\{\mathrm{abstain}\}$ may answer $a$ only when every candidate in the realized fiber $r^{-1}(x)$ has question value $a$. The maximal answerable domain is exactly the union of the fibers on which $q$ is constant. On that domain the answer is unique. Under a full-support prior this decoder uniquely maximizes the probability of answering, up to unrealized outputs.
\end{thmplain}
\begin{proof}
If a fiber contains two question values, any non-abstaining answer is wrong for at least one candidate. If the fiber is constant, its common value is safe. Thus decisions can be made independently on each realized fiber, and including every constant fiber gives the maximal domain. Full support makes omission of any such fiber strictly reduce retained probability.
\end{proof}

This adds a useful intermediate notion between total zero-error answerability and probabilistic loss: how much of the domain permits a certified answer. With non-full-support priors the maximal domain is unchanged, but probability-maximizing decoders need not be unique on zero-mass fibers.

\textbf{Semantic caution.} Equality of registered pieces can encode distinctions beyond what an individual piece entails. For $D=\{0,1\}$, registrations $\kappa(0)=\{0\}$, $\kappa(1)=D$ are sound and distinct. Observing which piece was emitted identifies the candidate, although the constraint $D$ alone entails no answer. The manuscript's equality-kernel answerability concerns observation of the registered representation. Constraint entailment is a different semantics and should be named when intended.

\subsecn{2}{D.2}{Saturation algebra and its support boundary}

\begin{thmplain}[\hypertarget{stmt:D.2}{Lemma~D.2}]
For a partition $\pi$, saturation satisfies
\[
s_\pi(\varnothing)=\varnothing,\qquad E\subseteq s_\pi(E),\qquad s_\pi\bigl(s_\pi(E)\cap F\bigr)=s_\pi(E)\cap s_\pi(F).
\]
If $\rho\le\pi$, then $s_\rho s_\pi=s_\rho=s_\pi s_\rho$. If $E$ is $\pi$-saturated and $F$ is $\rho$-saturated, then $E\cap F$ is $(\pi\vee\rho)$-saturated.
\end{thmplain}
\begin{proof}
A $\pi$-block meets $s_\pi(E)\cap F$ exactly when it meets both $E$ and $F$: once it meets $E$, the whole block lies in $s_\pi(E)$. This proves the displayed identity block by block. Nested partitions give nested block unions, proving both compositions. A $(\pi\vee\rho)$-block lies within one block of each partition, so membership in $E\cap F$ is constant on it.
\end{proof}

These identities support combination by intersection, units $D$, nulls $\varnothing$, and downward focusing by saturation on any join-closed domain family, in the generalized set-algebra sense. They do not require commuting extraction for incomparable domains.

\begin{thmplain}[\hypertarget{stmt:D.2a}{Proposition~D.2a}]
If the domain family is finite, contains a support of $E$, and is closed under ambient nonempty meets, then $E$ has a least support in that family. Join closure alone does not suffice.
\end{thmplain}
\begin{proof}
Saturation of $E$ under a partition says its equivalence relation never crosses from $E$ to its complement. The relation of the ambient meet is generated by unions of the supporting equivalence relations; every path still stays on one side. Thus $E$ is saturated for the meet of all its supports. Meet closure places that least support in the family. For a counterexample take $D=\{1,2,3,4\}$, $a=12|3|4$, $b=1|23|4$, and the join-closed label family $\{\bot,a,b,\top\}$. The piece $\{4\}$ is supported by $a,b,\top$ but not $\bot$, so has no least support. Its missing ambient-meet support is $a\wedge b=123|4$.
\end{proof}

\subsecn{3}{D.3}{Exactly when independent replication is idempotent}

\begin{thmplain}[\hypertarget{stmt:D.3}{Theorem~D.3}]
For a finite experiment $P$, the following are equivalent:
\begin{enumerate}
\item $P\otimes P$ is Blackwell equivalent to $P$, where the two outputs are independent given the state.
\item Every pair of rows of $P$ is either identical or has disjoint support.
\item The zero-error partition $\sigma^0(P)$ equals the equal-law partition $\sigma^{=}(P)$.
\item $P$ is Blackwell equivalent to the deterministic experiment reporting its row-equivalence class.
\end{enumerate}
\end{thmplain}
\begin{proof}
Give the state a full-support prior, and let $Y,Y'$ be the independent replicates. Under (1), data processing in both directions gives $I(Z;Y,Y')=I(Z;Y)$, so $I(Z;Y'\mid Y)=0$. For any realized $y$ and any state $z$ with $P(y\mid z)>0$, conditional independence of replicates gives $\mathrm{Law}(Y'\mid Z=z,Y=y)=P(\cdot\mid z)$. Vanishing conditional mutual information forces this law to depend only on $y$. Therefore states sharing any positive-probability output have identical rows, proving (2).

Under (2), each realized output identifies the row class, and sampling the common row from that class simulates $P$. This proves (4). The support-overlap graph then has exactly the row classes as components, giving (3). Conversely (3) forbids an overlap edge between unequal rows, proving (2). Finally (4) makes replication equivalent to repetition of a deterministic class label, proving (1).
\end{proof}

Thus strict accumulation under independent replication is exactly the failure of the experiment to be equivalent to deterministic information. This supplies a structural replacement for the erroneous ``posterior $1/2$ is impossible'' argument: a BSC can always be garbled to posterior $1/2$ by discarding its output. For crossover $c\in(0,1/2)$, two concordant outputs instead yield a posterior outside $[c,1-c]$, which one BSC cannot produce. At $c=0.3$, the independently checked deficiency is $0.084$.

\subsecn{4}{D.4}{Quotient averaging enlarges the feasible coupling choices}

Let $\mu$ have full support on $D$, and let $q:D\to A$ be surjective. For a fine-state channel $P$, define $(M_qP)(\cdot\mid a)=\sum_{z:\,q(z)=a}\mu(z\mid a)\,P(\cdot\mid z)$. Let $F_D$ be the fiber of joint channels with specified fine-state marginals, and $F_A$ the fiber with their averaged marginals. Write
\[
m_D=\min_{Q\in F_D} I_\mu\bigl(q(Z);Y\bigr),\qquad
m_A=\min_{R\in F_A} I_{\mu q}(A;Y).
\]

\begin{thmplain}[\hypertarget{stmt:D.4}{Theorem~D.4}]
$m_A\le m_D$. For an actual fine coupling $P$, its coarse fiber excess is at least its fine fiber excess, and their difference is exactly $m_D-m_A$. Equality holds precisely when some coarse minimizer has a lift in $F_D$.
\end{thmplain}
\begin{proof}
Averaging maps $F_D$ into $F_A$ and leaves the joint law of $(q(Z),Y)$ unchanged. Hence minimizing on its image cannot improve on minimizing over all $F_A$. Both fibers are compact finite-dimensional polytopes, and mutual information is continuous, so minima exist. Equality holds exactly when a fine minimizer maps to a coarse minimizer, equivalently when some coarse minimizer has a feasible lift. Subtract the two minima from the same actual mutual information to obtain the excess formula.
\end{proof}

\begin{thmplain}[\hypertarget{stmt:D.4a}{Lemma~D.4a}]
$\delta(M_qE,M_qF)\le\delta(E,F)$.
\end{thmplain}
\begin{proof}
Any fixed output simulator for the fine experiments also simulates their averages. Convexity of total variation bounds each averaged error by the maximum fine-state error. Infimize over simulators.
\end{proof}

Averaging therefore contracts simulation error, but need not preserve determinism or conditional independence. In the XOR example deterministic fine-state marginals force a unique fine coupling, while their question-averages admit different couplings. The excess difference is a concrete measure of a lost lifting constraint, not evidence that ordinary question-level information decompositions are undefined.

\subsecn{5}{D.5}{Least coupling and the Blackwell join}

\begin{thmplain}[\hypertarget{stmt:D.5}{Theorem~D.5}]
Fix finitely many finite marginal experiments $P_i$ on the same state space. Their full, unconstrained coupling fiber has a Blackwell-least member if and only if the $P_i$ have a least upper bound in the finite-experiment Blackwell order. When they exist, these two experiments are equivalent.
\end{thmplain}
\begin{proof}
Suppose a join $H$ exists. Choose simulators $G_i$ with $P_i=G_iH$, and, conditional on an output of $H$, sample the outputs of the $G_i$ independently. The resulting joint $Q$ lies in the coupling fiber and is a garbling of $H$. Its marginals show that it is an upper bound of all $P_i$, so $H\preceq_B Q$ as well. Every fiber member is an upper bound, hence dominates $H$ and therefore $Q$.

Conversely, suppose the fiber has a least member $L$. Any upper bound $H$ of the marginals yields a joint $Q$ by the same conditional simulation construction, with $Q\preceq_B H$. Since $L\preceq_B Q$, we have $L\preceq_B H$. As $L$ itself has the required marginals, it is their least upper bound.
\end{proof}

Consequently binary-state finite coupling fibers always have a least Blackwell class, by the binary lattice theorem \autocite[Proposition~16]{bertschinger2014}. The state-conditionally-independent coupling need not represent it. For two equal noisy marginals, the diagonal coupling is equivalent to one copy, whereas independent replication can be strictly stronger by D.3. Two incomparable fiber points do not exclude a least point elsewhere in the fiber. This theorem concerns the entire unconstrained fiber; budget, causal, or restricted-corruption subsets need a separate existence argument.

\subsecn{6}{D.6}{A sufficient transfer rule for graded budgets}

\begin{thmplain}[\hypertarget{stmt:D.6}{Proposition~D.6}]
Let the honest family satisfy graded separable bracketing under a joint ceiling $\Gamma_S$. Suppose a perceived family is obtained by a candidate-independent random choice $r$ followed by a product of local output channels $M_i^r$. Fix perceived local decoders $G_i$. If, for every branch of positive probability, each composite $G_iM_i^r$ is admissible under its honest singleton ceiling, then the joint perceived reading is a garbling of $\Gamma_S$.
\end{thmplain}
\begin{proof}
In each branch, honest separable bracketing supplies a simulator $L_r$ from $\Gamma_S$ to the joint decoded experiment. Mixing these simulators with the candidate-independent branch probabilities gives a simulator for the perceived joint reading.
\end{proof}

For a single product branch, admissibility of the perceived readings is exactly the needed composite admissibility. For a mixture, admissibility after averaging branches does not imply branchwise admissibility; it cannot be silently substituted. This is a useful sufficient criterion, not a characterization of all robust corruptions.

\textbf{Order caution.} The down-set completion embeds any poset by $x\mapsto{\downarrow}x$ and supplies unions as joins. It does not generally preserve joins already present: for incomparable $a,b$ with a join, ${\downarrow}a\cup{\downarrow}b$ omits $a\vee b$. In particular it encodes available alternatives, not an assembled joint experiment. The manuscript's $\min_i\delta(E_i,F)$ similarly chooses one local resource; it is not deficiency from a bundle of all resources.

\subsecn{7}{D.7}{Sharp committed-bit detection, including adaptive forgers}

\begin{thmplain}[\hypertarget{stmt:D.7}{Theorem~D.7}]
In each of $n$ rounds the honest audit reveals a fresh independent fair bit $Y_i$ and an honest verdict $A_i=Y_i$. Under forgery, $A_i$ must be chosen before $Y_i$ is revealed and may depend on the entire past and private randomness, but that information is independent of the fresh bit. For equal prior odds between honesty and forgery, the minimax error probability over all such adaptive forgers is exactly $2^{-n-1}$.
\end{thmplain}
\begin{proof}
Accept honesty if and only if every verdict matches. Honesty is never rejected. Conditional on any past and precommitted verdict, the fresh bit matches with probability $1/2$; iterated conditioning gives probability $2^{-n}$ of all matches under every allowed forgery. The equal-prior risk is therefore $2^{-n-1}$.

For a lower bound, allow the forger to choose independent fair verdict bits. Its joint law is uniform on all $4^n$ bit/verdict sequences. The honest law is uniform on the $2^n$ matching sequences and zero elsewhere. The optimal simple-hypothesis equal-prior Bayes error is one half the sum of the smaller probability at each outcome, namely $\tfrac12\,2^n\,4^{-n}=2^{-n-1}$. No test can beat that against this allowed forger.
\end{proof}

At $n=20$, the exact risk is $2^{-21}\approx4.77\times10^{-7}$, sharper than the original Bhattacharyya bound. This does not cover partial-round attacks, leaked fresh bits, biased/dependent future bits, or noisy honest verdicts. Those are separate models with potentially different rates.
